% !TEX root = main.tex
% 由文献综述_重写版_5000字.md转换；仅转换标记与排版，正文及引用保持原文。
\begin{center}
{\zihao{3}\bfseries 基于机器视觉的金属齿轮外观缺陷检测技术综述\par}
\end{center}

\textbf{摘要}：齿轮是机械设备中核心的传动零件，其表面的质量直接关系到传动的平稳性、噪音与使用寿命。在齿轮制造过程中，微小划痕与局部齿部缺损是两类经常出现的缺陷。过去主要依靠人工目视抽检，容易因疲劳产生漏检，且不同人员标准不一，难以满足现代工厂自动化全检的要求。本文针对金属零件与齿轮外观缺陷的机器视觉检测技术进行文献调研与综述。首先介绍了齿轮质检背景以及金属反光、加工刀纹和立体曲面等实际困难；然后简要介绍了传统的图像处理方法、深度学习检测方法以及无监督异常检测方法的基本原理与特点；接着针对金属反光、微小划痕识别、旋转多角度拍照判定以及缺陷样本少这四个核心问题，详细介绍了多篇代表性文献的具体实验、数据结果和结论；最后总结了各类方法的优缺点，并结合我们团队开展的 GearPro 齿轮质检项目，总结了对实际系统开发的四点工程启发。

\textbf{关键词}：齿轮外观检测；机器视觉；深度学习；目标检测；多视角判定；缺陷分类


\section*{一、 齿轮外观缺陷检测的研究背景}

\subsection*{1. 齿轮质检的意义与人工目检的问题}
在汽车、变速箱以及各种机械装备制造中，齿轮是传递动力和运动最基础的零件。齿轮在切削加工（滚齿、磨齿）、热处理和车间搬运时，表面容易产生各种外观缺陷。其中，表面浅划痕和局部的齿部缺损（如齿侧碰伤、齿角崩缺）是本文关注的两类缺陷。

如果带有这些缺陷的齿轮被装进机器，在高速旋转和交变冲击下，缺损边缘极易产生应力集中，不仅会让齿轮运转时的噪音变大、磨损变快，严重时还可能导致断齿损坏，造成安全隐患。

过去车间主要依靠质检工人用肉眼或放大镜进行人工抽检。但人工目视检查存在明显缺点：第一，工人长时间工作容易视觉疲劳，工作后期的漏检率明显增加；第二，每个人对划痕和磕碰的判定尺度不一，主观性较强；第三，抽检无法做到每个零件都检查，依然会有坏件流入下道工序。因此，开发自动化的机器视觉检测系统，用工业相机和算法代替人眼实现每件必检，是现代制造质量控制的重要方向[\hyperlink{ref:1}{1},\hyperlink{ref:2}{2}]。

\subsection*{2. 齿轮视觉检测遇到的实际困难}
将机器视觉用到金属齿轮外观检测时，会遇到几个很棘手的实际困难[\hyperlink{ref:3}{3}]：
\begin{itemize}
\item \textbf{金属反光和加工刀痕的干扰}：齿轮表面经车磨精加工后非常光滑，容易产生镜面反光。普通光照下易出现局部高光亮斑使浅划痕湮没；同时，磨削留下的细密刀纹在图像上与浅划痕很像，模型容易把正常刀纹误判为划痕。
\item \textbf{渐开线齿面曲率与齿根阴影}：齿轮表面为弯曲的渐开线齿廓，光线在不同位置的反射角度急剧变化；特别是狭窄的齿根过渡圆角处容易形成阴影遮挡与多次互反射，给均匀打光带来极大困难。
\item \textbf{需要旋转多角度拍摄与连续判断}：单个固定相机无法拍到立体齿轮的全貌，必须用旋转机构带动齿轮自转并配合多相机连续拍照[\hyperlink{ref:5}{5}]。这导致缺陷在多张图像中动态移动，算法不能只看单张照片，必须综合多张照片来决定整个齿轮是否合格。

\end{itemize}
\subsection*{3. 本文综述的主要内容与范围}
一套完整的视觉检测系统包括打光照明、相机拍照、图像算法处理以及机械电气分选等环节[\hyperlink{ref:1}{1}]。本文主要关注机械制造过程中，\textbf{怎样利用图像算法自动识别齿轮表面的微小划痕和齿部缺损，以及怎样根据多角度照片综合判断齿轮是否合格}。

需要说明的是，出厂验收涉及的齿形齿向公差测量，以及用超声波探伤内部缺陷、测量传动运转时的振动噪音和寿命预测等，不属于本文讨论的制造外观质检范围。

本文围绕“怎么拍清楚、怎么用算法识别出来，到怎么综合判断整个齿轮”的过程，先梳理常见检测方法，再精读关键文献，最后结合我们正在做的齿轮检测项目总结经验与启发。


\section*{二、 常见的表面缺陷检测方法介绍}

通过对文献的调研，目前的检测方法主要分为传统图像处理方法、深度学习监督检测方法以及无监督异常检测方法三大类[\hyperlink{ref:2}{2}]。

\subsection*{1. 传统的图像处理与人工特征方法}
在深度学习普及前，工业视觉主要依靠工程师手动设计特征提取公式和判断规则。通常先提取图像局部的边缘、灰度或纹理特征，再通过设定固定阈值或简单分类器（如支持向量机 SVM、极限学习机 ELM）完成判断。

常用方法包括提取微观纹理的局部二值模式（LBP）[\hyperlink{ref:6}{6},\hyperlink{ref:7}{7}]、灰度共生矩阵，以及提取局部角点特征的 SIFT 和 SURF 算子[\hyperlink{ref:7}{7}]。传统方法的优点是\textbf{原理简单直观、计算量小、运行速度非常快（通常几毫秒至十几毫秒）}，在打光完全固定、工件位置固定且缺陷规律的场景下表现良好。但其缺点是对光照变化极度敏感，车间光照稍微变化或零件稍有偏转就容易失效，换零件时重新调参非常繁琐。

\subsection*{2. 深度学习监督检测方法}
深度卷积神经网络（CNN）出现后，不再需要人工编写提取规则，而是由网络自动从大量样本中学习特征，对光照和复杂背景的适应能力大大增强[\hyperlink{ref:2}{2}]。根据输出形式不同，主要有三种任务：
\begin{itemize}
\item \textbf{图像分类（如 ResNet、EfficientNet）}：输出整张图片有无缺陷的概率[\hyperlink{ref:8}{8},\hyperlink{ref:11}{11},\hyperlink{ref:12}{12}]。优点是速度快、模型小且标注最省事（只需打合格/不合格标签）；缺点是无法给出缺陷具体位置，容易漏掉细小划痕，适合产线快速粗筛。
\item \textbf{目标检测（如 Faster R-CNN、YOLO 系列）}：用矩形框把缺陷圈出并给出类别和置信度[\hyperlink{ref:13}{13},\hyperlink{ref:14}{14},\hyperlink{ref:21}{21}]。目标检测非常直观，方便人工复核，标注只需画框，是工业界使用最普遍的方案；但对细长弯曲划痕，矩形框会包含较多正常背景。
\item \textbf{语义分割（如 U-Net）}：对每个像素分类，精确描绘不规则缺损轮廓[\hyperlink{ref:15}{15}]。优点是能精确测量缺损面积；缺点是训练需要逐像素勾线抠图，标注成本极高，在多规格产线上很难推广。

\end{itemize}
监督学习效果好，但前提是\textbf{必须先收集到大量标好的真实缺陷图片}，而在实际工厂里获取这些缺陷样本往往代价很大。

\subsection*{3. 无监督异常检测方法}
针对工厂里“正常合格品占绝大多数（>99\%）、真实缺陷样本极难收集”的痛点，研究人员提出了无监督异常检测方法[\hyperlink{ref:9}{9}]。

其基本思路是：\textbf{不再学习各类缺陷，而是只学习正常工件的外观分布}。训练时只使用易收集的正常合格品照片建立特征基准；测试时对比新照片与正常模式的差异，偏离了正常规律就判定为异常并输出热力图。代表算法如 PaDiM[\hyperlink{ref:17}{17}]和 PatchCore[\hyperlink{ref:10}{10}]，PaDiM 利用预训练卷积网络提取局部图像块特征，在特征图的每个空间位置独立拟合一个多元高斯分布，测试时通过计算待测特征与正常高斯分布的马氏距离得到异常得分，既保留了空间位置信息，又省去了网络训练过程。PatchCore 利用预训练模型提取正常图片的局部特征构建特征记忆库，采用贪心核心集采样压缩存储空间；测试时通过计算待测特征与记忆库中邻近特征的最近邻欧氏距离来判定异常。

无监督异常检测省去了收集缺陷样本与人工画框标注的麻烦，能有效检测未知新缺陷；但难点在于容易受到工件正常装夹微小位移和正常刀纹波动的干扰，在复杂金属加工表面上容易产生误报警。


\section*{三、 齿轮外观检测的关键问题与代表性文献研究}

在实际齿轮检测中，往往面临金属反光、微小划痕拍不清、旋转多角度拍照以及缺陷样本少这四个具体问题。下面结合代表性文献进行介绍和对比。

\subsection*{1. 解决金属表面反光与加工纹理的研究}
金属表面强反光和切削刀痕是引起模型误报和漏检的主要物理原因。单纯靠改动网络结构很难彻底解决，必须在拍照打光上下功夫。

Saiz等[\hyperlink{ref:4}{4}]专门针对高反光金属零件表面缺陷检测进行了研究。为了消除反光干扰，他们搭建了基于光度立体视觉（Photometric Stereo）的拍照系统：在固定相机周围安装多个不同方向的可控 LED 光源，快速交替点亮并拍摄同一零件在不同打光角度下的多张照片。接着根据反射模型计算出零件表面的高低起伏图和曲率图，并与原始灰度照片合并成多通道图片，送入其设计的自定义卷积分割网络（并将 U-Net[\hyperlink{ref:15}{15}] 与 DFANet 作为基线对比）。

\begin{figure}[!htbp]
\centering
\includegraphics[width=0.82\linewidth,height=0.27\textheight,keepaspectratio]{../figures/fig1_photometric_modality_ablation.png}
\caption*{图 1 光度立体多模态（RGB）与单纹理模态（Texture-Only）在三种金属材质下的缺陷分割性能消融对比（数据引自文献[\hyperlink{ref:4}{4}]）}
\end{figure}

对比实验表明：只用普通单张灰度照片时，反光亮斑边缘极易被误判成缺损；而加入表面高低起伏信息后，模型能清楚分辨出反光发白区域并没有物理凹陷，从而准确区分出真实缺损和反光亮斑[\hyperlink{ref:4}{4}]。Wu等[\hyperlink{ref:18}{18}]在 2025 年构建的复杂零件数据集中，也同样证明了多角度受控打光对消除反光、凸显微弱缺陷的核心作用。这说明：\textbf{解决金属反光的第一道防线是在硬件打光上}，把凹凸特征拍清楚能从根本上降低算法难度。

\subsection*{2. 检测微小划痕与多尺度缺陷的研究}
齿轮表面缺陷大小差别很大：有些磕碰面积较大，而有些浅划痕只有几个像素宽，在卷积网络多次下采样后极易被抹平丢失。

Liu等[\hyperlink{ref:8}{8}]针对钢材表面缺陷分类，提出轻量级多尺度卷积网络 ConCNN，该网络设计了 CNN-1 与 CNN-2 两路结构完全相同的卷积分支，分别输入 $200\times 200$ 与 $400\times 400$ 两种不同分辨率的图像，两路分支各自完成完整的卷积特征提取后，在全连接层前通过特征向量拼接实现晚期特征融合。在 NEU-CLS 六类钢材缺陷数据集（每类仅 60 张训练图）上的实验表明，ConCNN 分类准确率达到 98.89\%，相比未采用多尺度融合的单分支基线（92.14\%）提升了 6.75 个百分点[\hyperlink{ref:8}{8}]，证明在浅层保留高分辨率特征对识别细小缺陷非常有效。

Lin等[\hyperlink{ref:20}{20}]提出的特征金字塔（FPN）通过将深层特征传递回浅层特征图，让网络兼顾大缺陷与小划痕。Su等[\hyperlink{ref:19}{19}]在 2026 年提出的 MSFE-YOLO 将频域特征引入钢材检测，在 NEU-DET 上把 mAP@0.5 从 78.1\% 提升至 79.8\%；但频域计算使 RTX 3090 上的检测速度从 137.0 FPS 降至 89.3 FPS[\hyperlink{ref:19}{19}]。这告诉我们：在检测微小划痕时，\textbf{在 YOLO 中增设小目标 P2 检测层或采用多尺度特征融合能提升检出率}，但要控制模型复杂度以保证产线运行速度。

\begin{figure}[!htbp]
\centering
\includegraphics[width=0.9\linewidth]{../figures/fig2_msfe_yolo_pareto_tradeoff.png}
\caption*{图 2 YOLOv11s 基线与 MSFE-YOLO 在精确率、推理速度与参数量上的帕累托权衡关系（数据引自文献[\hyperlink{ref:19}{19}]）}
\end{figure}

\subsection*{3. 旋转多角度拍摄与齿轮整件判定的研究}
在专门针对齿轮立体检测的研究中，Allam等[\hyperlink{ref:5}{5}]针对汽车齿轮损伤搭建了包含旋转台和上下双相机的系统，在旋转中分别拍摄端面和侧面齿廓，并使用 Faster R-CNN 识别齿部损伤。

为解决单张照片因偶发反光引起的误报，作者提出了\textbf{连续多帧检出判定规则}——缺陷必须在连续旋转拍摄的多张照片中都被检出，才判定齿轮不合格。在 30 个缺陷齿轮和 100 个正常齿轮的实测中：
\begin{itemize}
\item 采用置信度阈值 0.90 且“连续三帧检出”规则时，成功检出全部 30/30 个缺陷齿轮（100\% 检出），但在 100 个正常齿轮中产生了 34 个误报；
\item 提高到“连续四帧检出”时误报减少，但漏检了部分仅在特定旋转角度短暂可见的轻微缺陷[\hyperlink{ref:5}{5}]。扫描单个齿轮耗时约 90 秒。

\end{itemize}
\begin{figure}[!htbp]
\centering
\includegraphics[width=0.9\linewidth,height=0.29\textheight,keepaspectratio]{../figures/fig3_temporal_voting_tradeoff.png}
\caption*{图 3 时序多帧投票机制下连续判定帧数对缺陷检出率与正常件误报数的双轴权衡曲线（数据引自文献[\hyperlink{ref:5}{5}]）}
\end{figure}

这给我们的启发是：\textbf{单张图片的算法检出率不等于整个齿轮的合格率}。多角度连续拍照利用多帧规则可过滤假报警；但规则太严又易漏检短暂可见缺陷；且 34\% 的误报率意味着每 100 个好零件就有 34 个被误拦，会给工人带来沉重复核负担。

\subsection*{4. 缺陷样本少与标注困难的解决方法}
针对手头样本少、微小划痕边界难以精确画准的问题，一些学者的实验给出了实用解决思路。

Tabernik等[\hyperlink{ref:16}{16}]在换向器裂纹检测中对比了两种标注方式（50 个工件，共 399 张图，其中仅 52 张有裂纹）：一种是逐像素描边的精细标注，另一种是直接画框的粗标注。实验结果表明：在最终判断“是否有裂纹”的任务中，用粗框训练出的模型效果与精细抠图标注几乎没有差别；同时训练时通过在训练迭代中交替抽取正样本与负样本实现跨迭代的平衡采样，有效稳定了小样本训练[\hyperlink{ref:16}{16}]。

Božič等[\hyperlink{ref:22}{22}]研究的混合监督方法表明，利用少量精确标注配合大量仅带“有/无缺陷”标签的图片联合训练，模型能以极低标注成本达到接近全精细标注的检测精度[\hyperlink{ref:22}{22}]。这指导我们：在样本有限且主要是筛出坏件时，\textbf{没必要花大量时间做像素级描边，采用粗框标注或暂时剔除连人眼都难界定的模糊样本（difficult 样本），同样能训练出稳定可靠的模型}。


\section*{四、 结论与对 GearPro 齿轮质检项目的启发}

\subsection*{1. 几种检测方法的对比与总结}
通过对上述文献的学习，我们将几种主要表面缺陷检测方法的优缺点与适用场景归纳在表 1 中。


\begin{table}[!htbp]
\centering
\caption*{\textbf{表 1　常见的工业表面视觉检测方法对比}}
\small
\renewcommand{\arraystretch}{1.2}
\setlength{\tabcolsep}{5pt}
\begin{tabularx}{\linewidth}{@{}>{\raggedright\arraybackslash}p{2.5cm}*{3}{>{\raggedright\arraybackslash}X}@{}}
\toprule
检测方法 & 特征与原理 & 训练标签需求 & 主要输出形式 \\
\midrule
\textbf{传统图像处理}[\hyperlink{ref:6}{6},\hyperlink{ref:7}{7}] & 灰度差分、LBP 纹理、手工阈值 & 样本很少；手工设定几何阈值 & 规则或二值判断 \\
\textbf{图像分类网络}[\hyperlink{ref:8}{8},\hyperlink{ref:11}{11}] & 卷积多层提取整图语义 & 图片级标签（合格/不合格） & 整图类别和置信度 \\
\textbf{目标检测网络}[\hyperlink{ref:5}{5},\hyperlink{ref:14}{14}] & 网格回归定位、FPN 特征金字塔 & 矩形框和缺陷类别 & 预测框位置、类别及置信度 \\
\textbf{无监督异常检测}[\hyperlink{ref:9}{9},\hyperlink{ref:10}{10},\hyperlink{ref:17}{17}] & 局部高斯分布、正常特征记忆库 & 仅需正常合格品图片 & 异常偏离分数和热力图 \\
\bottomrule
\end{tabularx}
\end{table}

\begin{table}[!htbp]
\centering
\caption*{\textbf{表 1（续）}}
\small
\renewcommand{\arraystretch}{1.2}
\setlength{\tabcolsep}{5pt}
\begin{tabularx}{\linewidth}{@{}>{\raggedright\arraybackslash}p{2.5cm}*{3}{>{\raggedright\arraybackslash}X}@{}}
\toprule
检测方法 & 优点 & 缺点 & 适用场景 \\
\midrule
\textbf{传统图像处理}[\hyperlink{ref:6}{6},\hyperlink{ref:7}{7}] & 原理简单，计算极快，算力要求小 & 对光照和反光极敏感，换零件需重新调参 & 光照稳定、缺陷形态固定的工位 \\
\textbf{图像分类网络}[\hyperlink{ref:8}{8},\hyperlink{ref:11}{11}] & 标注成本最低，速度极快，适合粗筛 & 无法给出具体位置，易漏微小划痕 & 产线高速过检、已知缺陷快速分选 \\
\textbf{目标检测网络}[\hyperlink{ref:5}{5},\hyperlink{ref:14}{14}] & 能标出具体位置，方便复核，标注适中 & 矩形框包含背景，细长划痕定位较粗 & 需指出缺陷位置指导分拣与人工复核 \\
\textbf{无监督异常检测}[\hyperlink{ref:9}{9},\hyperlink{ref:10}{10},\hyperlink{ref:17}{17}] & 不需要缺陷样本，能防范未知新缺陷 & 易受正常刀纹波动或微小位移干扰 & 缺陷极少、形态不固定的高精度初检 \\
\bottomrule
\end{tabularx}
\end{table}

\subsection*{2. 对我们做 GearPro 齿轮检测项目的启发}
通过学习上述文献的实验和结论，结合我们正在研发的 \textbf{GearPro 齿轮外观缺陷检测系统}（包含上下双相机、旋转机构、Jetson NX 工控机及算法），我们总结了以下四点实在的工程体会：

\begin{enumerate}
\item \textbf{先把光打好、把照片拍全，硬件是算法的基础}：
文献[\hyperlink{ref:4}{4}]和[\hyperlink{ref:18}{18}]表明，若照片反光严重或有死角没拍到，后面用多复杂的算法也很难补救。针对齿轮高反光、渐开线曲面和齿根阴影，必须采用漫射光源压制亮斑；同时借鉴 Allam 等[\hyperlink{ref:5}{5}]的双相机思路，我们在 GearPro 中采用上下两台工业相机配合自转，分别对准上端面、齿面和齿根，从硬件上尽量消除拍摄盲区。
\item \textbf{划痕和缺损分开建模，把分类和目标检测结合起来}：
划痕属细长弱对比缺陷，齿部缺损属块状几何缺陷。文献[\hyperlink{ref:8}{8}]和[\hyperlink{ref:20}{20}]启发我们，针对划痕检测在 YOLO 中增加保留高分辨率特征的 P2 小目标检测层；而在整体方案上，分类模型擅长大局判断（不易被局部噪点带偏），目标检测模型擅长找具体小框（不易漏小目标）。GearPro 把分类分数与 YOLO 检测分数加权融合，结合了两者的优点。
\item \textbf{不能只看单张图片，必须建立旋转连续多帧的判定规则}：
Allam 等[\hyperlink{ref:5}{5}]在齿轮上的实测说明，产线最终要判断的是整个齿轮能否使用，而非单张照片。齿轮旋转拍照时单张图片易因反光闪烁出现偶发误报。我们必须在软件中建立多帧滑动窗口规则，只有某个缺陷在旋转多张照片里连续被检出才触发报警，在保证不漏检的同时把误报压下来。
\item \textbf{测试集必须按齿轮独立划分，同时要严格控制误报率}：
Tabernik 等[\hyperlink{ref:16}{16}]证明剔除模糊标注框（difficult 框）有助于模型稳定训练；更重要的是，文献[\hyperlink{ref:23}{23}]提醒我们，测试模型时必须用未参与训练的新齿轮照片，绝对不能把同一个齿轮不同视角的照片混在训练和测试集里；文献[\hyperlink{ref:24}{24}]提醒我们，产线误报率过高会导致大量正常齿轮被误拦复核，因此挑选判定阈值时，必须在“尽量少漏检”和“控制工人复核工作量”之间找到合理平衡点。

\end{enumerate}

\begingroup
\setstretch{1.08}
\section*{参考文献}

{\small\noindent\hangindent=2em\hangafter=1 \hypertarget{ref:1}{}[1] NEWMAN T S, JAIN A K. A Survey of Automated Visual Inspection[J]. Computer Vision and Image Understanding, 1995, 61(2): 231–262.\par}

{\small\noindent\hangindent=2em\hangafter=1 \hypertarget{ref:2}{}[2] 罗东亮, 蔡雨萱, 杨子豪, 章哲彦, 周瑜, 白翔. 工业缺陷检测深度学习方法综述[J]. 中国科学: 信息科学, 2022, 52(6): 1002–1039.\par}

{\small\noindent\hangindent=2em\hangafter=1 \hypertarget{ref:3}{}[3] LUO Q, FANG X, LIU L, YANG C, SUN Y. Automated Visual Defect Detection for Flat Steel Surface: A Survey[J]. IEEE Transactions on Instrumentation and Measurement, 2020, 69(3): 626–644.\par}

{\small\noindent\hangindent=2em\hangafter=1 \hypertarget{ref:4}{}[4] SAIZ F A, BARANDIARAN I, ARBELAIZ A, GRAÑA M. Photometric Stereo-Based Defect Detection System for Steel Components Manufacturing Using a Deep Segmentation Network[J]. Sensors, 2022, 22(3): 882.\par}

{\small\noindent\hangindent=2em\hangafter=1 \hypertarget{ref:5}{}[5] ALLAM A, MOUSSA M, TARRY C, VERES M. Detecting Teeth Defects on Automotive Gears Using Deep Learning[J]. Sensors, 2021, 21(24): 8480.\par}

{\small\noindent\hangindent=2em\hangafter=1 \hypertarget{ref:6}{}[6] SONG K, YAN Y. A noise robust method based on completed local binary patterns for hot-rolled steel strip surface defects[J]. Applied Surface Science, 2013, 285: 858–864.\par}

{\small\noindent\hangindent=2em\hangafter=1 \hypertarget{ref:7}{}[7] LIU Y, XU K, WANG D. Online Surface Defect Identification of Cold Rolled Strips Based on Local Binary Pattern and Extreme Learning Machine[J]. Metals, 2018, 8(3): 197.\par}

{\small\noindent\hangindent=2em\hangafter=1 \hypertarget{ref:8}{}[8] LIU Y, YUAN Y, BALTA C, LIU J. A Light-Weight Deep-Learning Model with Multi-Scale Features for Steel Surface Defect Classification[J]. Materials, 2020, 13(20): 4629.\par}

{\small\noindent\hangindent=2em\hangafter=1 \hypertarget{ref:9}{}[9] BERGMANN P, FAUSER M, SATTLEGGER D, STEGER C. MVTec AD — A Comprehensive Real-World Dataset for Unsupervised Anomaly Detection[C]//Proceedings of the IEEE/CVF Conference on Computer Vision and Pattern Recognition (CVPR). 2019: 9592–9600.\par}

{\small\noindent\hangindent=2em\hangafter=1 \hypertarget{ref:10}{}[10] ROTH K, PEMULA L, ZEPEDA J, SCHÖLKOPF B, BROX T, GEHLER P. Towards Total Recall in Industrial Anomaly Detection[C]//Proceedings of the IEEE/CVF Conference on Computer Vision and Pattern Recognition (CVPR). 2022: 14318–14328.\par}

{\small\noindent\hangindent=2em\hangafter=1 \hypertarget{ref:11}{}[11] NAGY A M, CZÚNI L. Classification and Fast Few-Shot Learning of Steel Surface Defects with Randomized Network[J]. Applied Sciences, 2022, 12(8): 3967.\par}

{\small\noindent\hangindent=2em\hangafter=1 \hypertarget{ref:12}{}[12] ARIKAN S, VARANASI K, STRICKER D. Surface Defect Classification in Real-Time Using Convolutional Neural Networks[EB/OL]. arXiv:1904.04671, 2019.\par}

{\small\noindent\hangindent=2em\hangafter=1 \hypertarget{ref:13}{}[13] REN S, HE K, GIRSHICK R, SUN J. Faster R-CNN: Towards Real-Time Object Detection with Region Proposal Networks[C]//Advances in Neural Information Processing Systems (NeurIPS 28). 2015.\par}

{\small\noindent\hangindent=2em\hangafter=1 \hypertarget{ref:14}{}[14] REDMON J, DIVVALA S, GIRSHICK R, FARHADI A. You Only Look Once: Unified, Real-Time Object Detection[C]//Proceedings of the IEEE Conference on Computer Vision and Pattern Recognition (CVPR). 2016: 779–788.\par}

{\small\noindent\hangindent=2em\hangafter=1 \hypertarget{ref:15}{}[15] RONNEBERGER O, FISCHER P, BROX T. U-Net: Convolutional Networks for Biomedical Image Segmentation[C]//MICCAI 2015. 2015: 234–241.\par}

{\small\noindent\hangindent=2em\hangafter=1 \hypertarget{ref:16}{}[16] TABERNIK D, ŠELA S, SKVARČ J, SKOČAJ D. Segmentation-based deep-learning approach for surface-defect detection[J]. Journal of Intelligent Manufacturing, 2020, 31: 759–776.\par}

{\small\noindent\hangindent=2em\hangafter=1 \hypertarget{ref:17}{}[17] DEFARD T, SETKOV A, LOESCH A, AUDIGIER R. PaDiM: A Patch Distribution Modeling Framework for Anomaly Detection and Localization[C]//ICPR Workshops. 2021: 475–489.\par}

{\small\noindent\hangindent=2em\hangafter=1 \hypertarget{ref:18}{}[18] WU L, RAN Y, YAN L, LIU Y, SONG Y, HAN D. A dataset for surface defect detection on complex structured parts based on photometric stereo[J]. Scientific Data, 2025, 12: 276.\par}

{\small\noindent\hangindent=2em\hangafter=1 \hypertarget{ref:19}{}[19] SU S, SHEN J, LIN P, TANG W, ZHANG W, CHEN Z. MSFE-YOLO: A Steel Surface Defect Detection Algorithm Integrating Multi-Scale Frequency Domain and Defect-Aware Attention[J]. Sensors, 2026, 26(8): 2311.\par}

{\small\noindent\hangindent=2em\hangafter=1 \hypertarget{ref:20}{}[20] LIN T Y, DOLLÁR P, GIRSHICK R, HE K, HARIHARAN B, BELONGIE S. Feature Pyramid Networks for Object Detection[C]//Proceedings of the IEEE Conference on Computer Vision and Pattern Recognition (CVPR). 2017: 2117–2125.\par}

{\small\noindent\hangindent=2em\hangafter=1 \hypertarget{ref:21}{}[21] DING R, DAI L, LI G, LIU H. TDD-net: a tiny defect detection network for printed circuit boards[J]. CAAI Transactions on Intelligence Technology, 2019, 4(2): 110–116.\par}

{\small\noindent\hangindent=2em\hangafter=1 \hypertarget{ref:22}{}[22] BOŽIČ J, TABERNIK D, SKOČAJ D. Mixed supervision for surface-defect detection: From weakly to fully supervised learning[J]. Computers in Industry, 2021, 129: 103459.\par}

{\small\noindent\hangindent=2em\hangafter=1 \hypertarget{ref:23}{}[23] CAWLEY G C, TALBOT N L C. On Over-fitting in Model Selection and Subsequent Selection Bias in Performance Evaluation[J]. Journal of Machine Learning Research, 2010, 11: 2079–2107.\par}

{\small\noindent\hangindent=2em\hangafter=1 \hypertarget{ref:24}{}[24] SAITO T, REHMSMEIER M. The Precision-Recall Plot Is More Informative than the ROC Plot When Evaluating Binary Classifiers on Imbalanced Datasets[J]. PLOS ONE, 2015, 10(3): e0118432.\par}
\endgroup
